\section*{الفصل \ref{ch:g12:phylogenetic-trees} --- قراءة القرابة: الأشجار التطورية}

\begin{solution}{exo:g12:phylogenetic-trees:1}
الأطراف: \omterm{def:g10:biodiversity-scales:species}{الأنواع} المقارنة. والعقد: أسلافها المشتركة الأحدث المفترضة.
والجذر: السلف المشترك لها جميعًا.
\end{solution}

\begin{solution}{exo:g12:phylogenetic-trees:2}
أقرب أقرباء السحلية هما الفأر والإنسان معًا (فهما يشتركان في العقدة
التي تقع تحت فرع السحلية مباشرة). وأقرب أقرباء الضفدع مجموعة \omterm{def:g10:biodiversity-scales:species}{الأنواع}
الثلاثة الأخرى كلها، على السواء.
\end{solution}

\begin{solution}{exo:g12:phylogenetic-trees:3}
حالة ظهرت عند سلف وورثها نسله؛ فالاشتراك فيها اشتراك في ذلك السلف. أما
الحالة السلفية فكانت عند سلف المجموعة كلها، فكل \omterm{def:g10:biodiversity-scales:species}{نوع} كان يستطيع أن
يبقيها: والاشتراك فيها لا يبين شيئًا عن قرابة أوثق.
\end{solution}

\begin{solution}{exo:g12:phylogenetic-trees:4}
سلف وكل نسله. والطيور \omterm{prop:g12:phylogenetic-trees:rules}{فرع حيوي}. أما "الزواحف" بالمعنى الدارج فليست
كذلك: فهي تترك خارجها الطيور، وهي منحدرة من داخلها.
\end{solution}

\begin{solution}{exo:g12:phylogenetic-trees:5}
الشمبانزي؛ وأقدم العقد هي التي تضم المكاك إلى القردة العليا، قبل 25
مليون سنة.
\end{solution}

\begin{solution}{exo:g12:phylogenetic-trees:6}
أدر عند كل عقدة: فينشق الجذر إلى الضفدع (في الأسفل) وإلى الباقي؛ ثم
السحلية، ثم الفأر والإنسان في الأعلى. ويظل الفأر والإنسان يشتركان في
أحدث عقدة، وتنضم إليهما السحلية بعد ذلك، ثم الضفدع آخرًا.
\end{solution}

\begin{solution}{exo:g12:phylogenetic-trees:7}
للحمامة فكوك وأطراف وسلى: فهي تدخل \omterm{prop:g12:phylogenetic-trees:rules}{الفرع الحيوي} للسلويات إلى جانب
السحلية والفأر. أما الريش، وهو عند \omterm{def:g10:biodiversity-scales:species}{نوع} واحد، فلا يحدد مجموعة هنا؛ وكان
سيحدد الطيور لو أضيف مزيد من الطيور.
\end{solution}

\begin{solution}{exo:g12:phylogenetic-trees:8}
الشعر واللبن والفك تضع الدلفين بين الثدييات؛ ووضعه مع القرش كان
سيقتضي أن تكون الصفات الثديية كلها قد نشأت مرتين. أما الزعانف والجسم
الانسيابي فنشآ على حدة في السلالتين: أي إنهما تقارب.
\end{solution}

\begin{solution}{exo:g12:phylogenetic-trees:9}
$2.2/0.27 \approx 8$ ملايين سنة.
\end{solution}

\begin{solution}{exo:g12:phylogenetic-trees:10}
كلاهما طرف حاضر؛ ولا واحد منهما منحدر من الآخر. والصواب: الإنسان
والشمبانزي نوعان شقيقان، منحدران من سلف مشترك، انقرض الآن، عاش قبل
نحو ستة ملايين سنة ولم يكن إنسانًا ولا شمبانزي.
\end{solution}

\begin{solution}{exo:g12:phylogenetic-trees:11}
الجلكى يعلم أنه يقع خارج مجموعة \omterm{def:g10:common-ancestry:bodyplan}{الفقاريات} ذوات الفكوك، فحالاته سلفية
لها. ولو اتخذ الفأر مجموعةً خارجية لقرئ الشعر والسلى والأطراف والفكوك
كلها سلفية وقرئ غيابها مشتقًا: فكانت الشجرة ستبنى مقلوبة.
\end{solution}

\begin{solution}{exo:g12:phylogenetic-trees:12}
كانت \omterm{def:g12:selection-drift-speciation:population}{الجمهرة} السلفية تحمل عدة \omterm{def:g10:universal-dna:gene}{أليلات} لكل \omterm{def:g10:universal-dna:gene}{مورّثة}؛ فلما انشقت مرتين في
تعاقب سريع، فرزت \omterm{def:g10:universal-dna:gene}{أليلات} مختلفة إلى \omterm{def:g10:biodiversity-scales:species}{الأنواع} المنحدرة الثلاثة عند
\omterm{def:g10:universal-dna:gene}{مورّثات} مختلفة، فصار تاريخ \omterm{def:g10:universal-dna:gene}{مورّثة} يجمع \omterm{def:g10:biodiversity-scales:species}{النوع} 1 مع \omterm{def:g10:biodiversity-scales:species}{النوع} 2 وتاريخ أخرى
يجمع 2 مع 3. وكل شجرة \omterm{def:g10:universal-dna:gene}{مورّثة} صحيحة لتلك \omterm{def:g10:universal-dna:gene}{المورّثة}؛ أما شجرة \omterm{def:g10:biodiversity-scales:species}{الأنواع}
فهي حكم الأغلبية من \omterm{def:g10:universal-dna:gene}{مورّثات} كثيرة.
\end{solution}

\begin{solution}{exo:g12:phylogenetic-trees:13}
عند 1\% في السنة، تختلف عينات فيروسية أخذت بفارق شهور اختلافًا قابلًا
للقياس أصلًا: فشجرة الفيروس تفصل في من أعدى من داخل وباء، لكن التسلسل
كان سيمحى ويكتب مرات كثيرة على مدى ملايين السنين. وعند 0.3\% لكل مليون
سنة، تتغير \omterm{def:g10:universal-dna:gene}{مورّثة} ثديية أبطأ من أن يرى شيء داخل وباء، لكنها تحفظ سجلًا
مقروءًا على مدى 200 مليون سنة.
\end{solution}

\begin{solution}{exo:g12:phylogenetic-trees:14}
الأحفورة \omterm{def:g10:biodiversity-scales:species}{نوع} قائم بذاته، له حالاته المشتقة، وليس بالضرورة \omterm{def:g12:selection-drift-speciation:population}{الجمهرة}
السلفية؛ وهي تقع على فرع جانبي قرب العقدة التي ظهر عندها الريش. وهي
تبين أن العقدة --- أي السلف المشترك للطيور وأقرب أقربائها من الزواحف
--- عمرها 150 مليون سنة على الأقل، وأن الريش سبق فقد الأسنان والذنب.
\end{solution}

\begin{solution}{exo:g12:phylogenetic-trees:15}
لا قمة للشجرة: فإدارة أي عقدة تنقل الإنسان إلى الأسفل من غير أن تغير
قرابة واحدة. وكل طرف حاضر تطور المدة نفسها منذ الجذر، فلا واحد منها
أرقى. والإنسان طرف واحد بين ملايين، على فرع واحد من القردة العليا؛
وشكل الشجرة يسجل الانحدار لا الرتبة.
\end{solution}

\begin{solution}{pb:g12:phylogenetic-trees:1}
\textbf{1.} الهيكل العظمي (خمسة \omterm{def:g10:biodiversity-scales:species}{أنواع}): أي \omterm{def:g10:common-ancestry:bodyplan}{الفقاريات} العظمية؛ والقرش
خارجها.

\textbf{2.} أربعة أطراف: رباعيات الأطراف (السمندل والتمساح والحمامة
والقط). والبيضة السلوية: السلويات (التمساح والحمامة والقط). والتداخل:
\omterm{def:g10:common-ancestry:bodyplan}{الفقاريات} العظمية $\supset$ رباعيات الأطراف $\supset$ السلويات.

\textbf{3.} لا: فحالة عند \omterm{def:g10:biodiversity-scales:species}{نوع} واحد لا تجمع شيئًا. وكانت ستفيد لو
اشتركت فيها \omterm{def:g10:biodiversity-scales:species}{أنواع} أكثر (طيور أخرى، وثدييات أخرى، وزواحف أخرى).

\textbf{4.} الجذر: القرش في مقابل الباقي (الهيكل العظمي)؛ ثم السلمون
في مقابل رباعيات الأطراف (أربعة أطراف)؛ ثم السمندل في مقابل السلويات
(البيضة السلوية)؛ وبين السلويات انشقاق ثلاثي غير محسوم، مع الريش على
فرع الحمامة والشعر والفتحة الجمجمية على فرع القط.

\textbf{5.} الجدول يضع التمساح والحمامة والقط معًا لكنه لا يستطيع
ترتيبها: فأقرب أقرباء التمساح غير محسوم.

\textbf{6.} التمساح--الحمامة، 8\%.

\textbf{7.} القط (15\% من كل منهما)؛ ثم السمندل (19--20\%)؛ ثم
السلمون (24--25\%)؛ ثم القرش (30--31\%).

\textbf{8.} تتفقان في كل تداخل؛ ويفصل \omterm{def:g10:universal-dna:information}{DNA} في السلويات: فالتمساح
والحمامة شقيقان، والقط ينضم إليهما بعدهما.

\textbf{9.} انفصلت سلالة القرش قبل أن تتباعد سائر السلالات بعضها عن
بعض؛ فكل منها بقي مفارقًا للقرش المدة نفسها ويبدي المسافة نفسها.

\textbf{10.} التمساح أقرب قرابة إلى الحمامة منه إلى القط --- وبالمنطق
نفسه منه إلى سحلية؛ وليست "الزواحف" من دون الطيور فرعًا حيويًا.

\textbf{11.} 19\% على مدى $2 \times 340$ مليون سنة من الفروع: أي نحو
0.028\% لكل مليون سنة لكل سلالة، أي 0.056\% لكل مليون سنة من
الانفصال.

\textbf{12.} التمساح--الحمامة $8/0.056 \approx 140$ مليون سنة؛ وعقدة
السلويات $15/0.056 \approx 270$ مليون سنة.

\textbf{13.} السلمون $24.5/0.056 \approx 440$؛ والقرش $30.5/0.056
\approx 540$ مليون سنة.

\textbf{14.} الأحافير: 250 و310؛ والساعة تعطي 140 و270. فالترتيب صحيح
والتاريخ الأقدم قريب، لكن تاريخ انفصال التمساح والطير أصغر بكثير مما
ينبغي: فمعدل \omterm{def:g10:universal-dna:gene}{مورّثة} واحدة يتغير بين السلالات، \omterm{def:g10:universal-dna:gene}{ومورّثة} واحدة تعطي في
أحسن الأحوال ساعة تقريبية.

\textbf{15.} عند المسافات الكبيرة تكون مواضع كثيرة قد تغيرت أكثر من
مرة، فلا تتجامع الفروق بعد: أي إن الساعة تتشبع. فالعقد القديمة تقدَّر
بأقل من قدرها بهذه \omterm{def:g10:universal-dna:gene}{المورّثة}؛ ويلزم \omterm{def:g10:universal-dna:gene}{مورّثة} أبطأ، أو \omterm{def:g10:universal-dna:gene}{مورّثات} كثيرة.

\textbf{16.} (القط، (التمساح، الحمامة)) هو أصغر \omterm{prop:g12:phylogenetic-trees:rules}{فرع حيوي} يحوي القط:
أي السلويات؛ ثم رباعيات الأطراف؛ ثم \omterm{def:g10:common-ancestry:bodyplan}{الفقاريات} العظمية؛ ثم الستة كلها
مع القرش.

\textbf{17.} "الأسماك" ليست فرعًا حيويًا: فالعقدة التي تضم القرش
والسلمون هي الجذر، ونسلها يشمل كل شيء. أما السلويات \omterm{prop:g12:phylogenetic-trees:rules}{ففرع حيوي}:
التمساح والحمامة والقط وسلفها المشترك.

\textbf{18.} السمندل \omterm{def:g10:biodiversity-scales:species}{نوع} حاضر لا وسط بين شيئين؛ وهو المجموعة الشقيقة
للسلويات، يشترك معها في سلف رباعيات الأطراف ويفتقر إلى البيضة السلوية
التي نشأت لاحقًا في خطها.

\textbf{19.} إلى جانب الفرع الواقع بين عقدة رباعيات الأطراف وعقدة
السلويات: أي رباعي أطراف ليس سلويًا. وتاريخها يؤكد أن الأطراف الأربعة
كانت موجودة قبل 360 مليون سنة وأن البيضة السلوية ظهرت لاحقًا.

\textbf{20.} (القرش، (السلمون، (السمندل، (القط، (التمساح،
الحمامة)))))؛ وقد فصل \omterm{def:g10:universal-dna:information}{DNA} في أن التمساح والحمامة شقيقان؛ ويؤرخ
انفصالهما، بهذه \omterm{def:g10:universal-dna:gene}{المورّثة}، بنحو 140 مليون سنة (وتقول الأحافير 250).
\end{solution}
